\begin{abstract}
	A comprehensive theoretical framework is developed for potential-field transformations 
	based on planar equivalent layers, a technique empirically used in applied geophysics since 
	the 1960s. The formulation replaces the widely accepted equivalent-layer principle -- primarily 
	founded on empirical generalizations -- with a rigorous theoretical foundation. The planar 
	equivalent-layer theorem establishes the existence of continuous planar layers capable of 
	reproducing potential-field data within a predefined error level, even when the layer 
	intersects the true sources. Closed-form expressions for planar equivalent layers are derived 
	by modifying the classical solutions of Dirichlet’s and Neumann’s problems, revealing that 
	these layers possess tangible physical meaning related to boundary-value problems in potential 
	theory. The theoretical framework also demonstrates the mathematical equivalence between 
	potential-field transformations computed via the equivalent-layer technique and those obtained 
	through Fourier-domain approaches. Distinct classes of equivalent layers are formally defined, 
	together with their relationships to the observed data, true sources, and potential-field 
	transformations. The resulting formulation provides a rigorous explanation for the 
	long-standing success of the equivalent-layer technique in potential-field analysis.
\end{abstract}